% Pursuing Stacks (cotes 134-2 to 134-8): archive cote 134-2, site page 65 (PDF page 66).
% appendix (letters to Breen). Draft compared: ps1.tex:1484-1527.
% AI-transcribed by Claude Opus 5.5 (claude-opus-5-5, adaptive thinking, effort high), Message Batches API, batch ps1 (msgbatch_01Rv5hCaNTtUt4JHx1yqoHPL), 2026-11-12; not yet checked by eye.
\scanpage{134-2/65}
\mspage{- 50 -}

\diagram{\begin{tikzcd}[row sep=small]
 & \mathrm{Rf}_!(\hat{\tilde{F}}) \arrow[dl] & \\
\mathrm{Rf}_!(F) \arrow[rr] & & \mathrm{R}\tilde{f}_*(\tilde{F}) \arrow[ul]
\end{tikzcd}}\ ,

where $\hat{\tilde{F}}$ is the \emph{formal completion} of $\tilde{F}$ along $\tilde{X} - X$ (a finite number of
points ...). It is here, unless I am mistaken, that appears the link with
\margin{bracket with an upward arrow from this line to the next}local class field theory, in its cohomological version, on which I am going
\margin{(espace)}\uline{now to say} \deleted{several}\inserted{a few} words. \qquad * \qquad * \qquad *

\margin{(App. 13)}\uline{\over{(E)}{\unsure{(App 13)}} \deleted{Local}\inserted{((E) Local} class field\deleted{s} \inserted{(theory} as a duality formula}

\margin{curved arrow from the margin note (App. 13) to this line; vertical double arrow from the heading line to the next line}Let $V$ be a complete discrete valuation ring \inserted{with} residue field $k$ - assume
either that $k$ \deleted{can be}\inserted{has been} lifted to $k \subset V$ (and therefore $V \simeq k[[T]]$) or
that $k$ is perfect\inserted{ of char. $p>0$}. In order to fix ideas, and to be sure that I'm on solid
ground, I consider at first on $K$ (= the field of fractions of $V$) \emph{finite}
coefficients $F$ (as on $X$ previously) and I consider \inserted{the objects} $H^1(K, F)$\hand{,} or
$R\Gamma_K(F)$ . The main work to be done consists in defining an adequate category
of coefficients over $k$ (perhaps the same one as in (D)) and a functor
\[ F \longmapsto \mathbb{R}\underline{\Gamma}_K(F) \]
with values in the category of such coefficients, in such a manner that the
following isomorphism holds.
\[ R\Gamma_K(F) \simeq R\Gamma_k(R\underline{\Gamma}_K(F)) \ . \]

This corresponds to the intuition (acquired directly from elementary
examples) according to which for $k$ algebraically closed, say, the
\margin{w/}$H^0(K, F)$ , $H^1(K, F)$ ... are endo\over{\unsure{ss}}{w}ed with a structure of $k$-algebraic group
(ind-pro ...). In this construction, the ring scheme of Witt vectors over
$k$ (introduced by Serre) and the ``Greenberg functor'' (associating to a
$V$-scheme a $k$-prescheme) will play an essential role.

This being done, the duality formula will be formally stated as in (D)
above:
\[ \Delta R\underline{\Gamma}_K(F) \simeq R\underline{\Gamma}_K(DF[1]) \]
where $D$ stands for Cartier duality, $\Delta$ for Serre duality. When the

\margin{65}
